% =============================================================================
%  CAPÍTULO 4 — RESULTADOS E DISCUSSÃO
%  TCC: Dollar Bars vs. Time Bars — Modelos GARCH para BTC/USDT
% =============================================================================

\chapter{Resultados e Discussão}
\label{cap:resultados}

% -----------------------------------------------------------------------------
\section{Propriedades Estatísticas Comparativas das Séries}
\label{sec:prop_estatisticas}

A análise exploratória das séries de retornos logarítmicos revela diferenças estruturais
substanciais entre as duas formas de amostragem, conforme sintetizado na \autoref{tab:descritivas}.
Essas diferenças não são ruído amostral: elas decorrem diretamente do mecanismo pelo qual
cada tipo de barra é construída.

\begin{table}[H]
\centering
\caption{Estatísticas descritivas dos retornos logarítmicos.}
\label{tab:descritivas}
\begin{tabular}{lSS}
\toprule
{Estatística} & {\emph{Dollar Bars}} & {\emph{Time Bars}} \\
\midrule
Média ($1^{\circ}$ momento)        &  3.46e-6  &  3.44e-6  \\
Variância ($2^{\circ}$ momento)    &  2.49e-6  &  2.40e-6  \\
Assimetria ($3^{\circ}$ momento)   &  0.0547   & -1.0071   \\
Curtose de Fisher ($4^{\circ}$ momento) &  3.0879   &  75.3687  \\
\bottomrule
\end{tabular}
\fonte{Elaboração própria.}
\end{table}

O indicador mais expressivo é a curtose. As time bars de 5 minutos apresentam curtose de Fisher
de 75,37, ao passo que as dollar bars registram apenas 3,09 — uma razão de aproximadamente
25 vezes. Para compreender essa disparidade, é preciso considerar o que ocorre durante
choques de volatilidade no mercado de Bitcoin. Ao longo de um episódio de alta atividade
financeira, dezenas de transações são comprimidas em uma única barra temporal, produzindo
retornos extremamente grandes que inflariam artificialmente as caudas da distribuição.
A dollar bar, ao encerrar uma barra sempre que o volume financeiro transacionado atinge o
limiar $\bar{D}$, redistribui esse mesmo choque ao longo de múltiplas barras, diluindo o
extremo. Esse mecanismo é precisamente a sub-amostragem em períodos de alta atividade
descrita por \citeonline{lopez2018advances} na discussão sobre barras de informação (seção 2.1.1).

A assimetria reforça o mesmo argumento. Enquanto as dollar bars exibem assimetria de apenas
0,05 — praticamente simétrica —, as time bars apresentam assimetria de $-1{,}01$. O sinal
negativo indica que quedas abruptas são capturadas de forma concentrada nessa amostragem:
grandes retornos negativos, correspondentes a correções rápidas, tendem a ocorrer em janelas
de poucos minutos e, por isso, se concentram em uma única barra temporal. Além da implicação
distribucional, a simetria das dollar bars facilita a especificação do modelo GARCH, que,
em sua formulação padrão com inovações gaussianas, assume distribuição condicional simétrica.

\begin{table}[H]
\centering
\caption{Testes de normalidade, estacionariedade e
heteroscedasticidade sobre os retornos logarítmicos.}
\label{tab:testes}
\begin{tabular}{lSSSS}
\toprule
{Teste} &
\multicolumn{2}{c}{\emph{Dollar Bars}} &
\multicolumn{2}{c}{\emph{Time Bars}} \\
\cmidrule(lr){2-3} \cmidrule(lr){4-5}
& {Estat.} & {\emph{p}-valor} & {Estat.} & {\emph{p}-valor} \\
\midrule
Jarque-Bera          & 83384.65  & {$<0{,}001$} & 49860565.37 & {$<0{,}001$} \\
Kolmogorov-Smirnov   & 0.0123    & {$<0{,}001$} & 0.0934      & {$<0{,}001$} \\
\addlinespace
ADF                  & -102.46   & {$<0{,}001$} & -57.32      & {$<0{,}001$} \\
KPSS                 & 0.2053    & 0.10         & 0.2171      & 0.10         \\
\addlinespace
Ljung-Box ($r_t$, lag 10)   & 50.06    & {$<0{,}001$} & 155.03   & {$<0{,}001$} \\
Ljung-Box ($r_t^2$, lag 10) & 79528.20 & {$<0{,}001$} & 28237.22 & {$<0{,}001$} \\
ARCH de Engle (lag 10)      & 30420.71 & {$<0{,}001$} & 24064.78 & {$<0{,}001$} \\
\bottomrule
\end{tabular}
\nota{H$_0$ por teste: Jarque-Bera e KS --- normalidade;
ADF --- raiz unitária; KPSS --- estacionariedade;
Ljung-Box e ARCH de Engle --- ausência de autocorrelação/efeito ARCH.}
\fonte{Elaboração própria.}
\end{table}

\begin{figure}[H]
    \centering
    \includegraphics[width=0.95\textwidth]{Dollar_Bars_5m_normalidade.png}
    \caption{Histograma com curva normal e QQ-Plot --- Dollar Bars (5 min).}
    \label{fig:normalidade_dollar}
\end{figure}

\begin{figure}[H]
    \centering
    \includegraphics[width=0.95\textwidth]{Time_Bars_5m_normalidade.png}
    \caption{Histograma com curva normal e QQ-Plot --- Time Bars (5 min).}
    \label{fig:normalidade_time}
\end{figure}

\begin{figure}[H]
    \centering
    \includegraphics[width=0.95\textwidth]{Dollar_Bars_5m_acf_pacf_retornos2.png}
    \caption{ACF e PACF dos retornos ao quadrado --- Dollar Bars (5 min).}
    \label{fig:acf_dollar_sq}
\end{figure}

\begin{figure}[H]
    \centering
    \includegraphics[width=0.95\textwidth]{Time_Bars_5m_acf_pacf_retornos2.png}
    \caption{ACF e PACF dos retornos ao quadrado --- Time Bars (5 min).}
    \label{fig:acf_time_sq}
\end{figure}

Os testes formais de adequação distribucional, apresentados na \autoref{tab:testes},
amplificam esse quadro. A estatística de Jarque-Bera das time bars é superior a 49 milhões,
contra aproximadamente 83 mil para as dollar bars — uma diferença de mais de 500 vezes.
Esse resultado não representa apenas rejeição da normalidade, que ambas as séries exibem:
indica que as caudas das time bars pertencem a uma ordem de magnitude distinta. O teste
de Ljung-Box nos retornos em nível rejeita a hipótese de ausência de autocorrelação
para as duas séries, mas com estatística consideravelmente maior para as time bars
(155,03 versus 50,06). Contudo, conforme discutido na seção seguinte, essa autocorrelação
é de magnitude modesta para o volume de observações disponível e não justifica componentes
autorregressivos ou de médias móveis no modelo final. Ambas as séries são estacionárias:
o teste Dickey-Fuller Aumentado rejeita a hipótese de raiz unitária, e o teste KPSS não
rejeita estacionariedade em nenhuma das duas amostras, validando o uso de $d = 0$ na
especificação ARMA.

Em conjunto, os resultados exploratórios corroboram a hipótese de \citeonline{lopez2018advances}:
dollar bars produzem séries com propriedades distribucionais superiores em termos de
curtose, assimetria e dependência linear residual. Cabe ressalvar, contudo, que propriedades
descritivas melhores não garantem, por si só, maior poder preditivo fora da amostra.
Essa questão é investigada nas seções seguintes.

% -----------------------------------------------------------------------------
\section{Seleção de Modelos via CPCV}
\label{sec:selecao_modelos}

O procedimento de seleção de hiperparâmetros via \textit{Combinatorial Purged Cross-Validation}
(CPCV) com $N = 10$ grupos, $K = 5$ grupos de teste e embargo de 240 barras gerou
$\binom{10}{5} = 252$ folds de avaliação, totalizando 36.288 ajustes MLE por série
(252 folds $\times$ 144 combinações de parâmetros no grid
$\text{ARMA}(0..5,\,0..5) \times \text{GARCH}(1..2,\,1..2)$).

\begin{table}[H]
\centering
\caption{Resultados da seleção de parâmetros via CPCV.}
\label{tab:cpcv_resultados}
\begin{tabular}{llrl}
\toprule
{Série} & {Modelo selecionado} & {Frequência} & {AIC médio} \\
\midrule
\emph{Dollar Bars} & ARMA(0,0)-GARCH(1,2) & 233/252 (92,5\%) &
    $-835.475$ \\
\emph{Time Bars}   & ARMA(0,0)-GARCH(1,1) & 252/252 (100\%) &
    $-859.083$ \\
\bottomrule
\end{tabular}
\fonte{Elaboração própria.}
\end{table}

\begin{figure}[H]
    \centering
    \includegraphics[width=0.95\textwidth]{comparacao_modelos/cpcv_01_aic_por_fold.png}
    \caption{AIC médio por fold CPCV --- Dollar Bars vs.\ Time Bars.}
    \label{fig:cpcv_aic}
\end{figure}

Em ambas as séries, o componente ARMA selecionado foi ARMA$(0,0)$, ou seja, média constante
sem termos autorregressivos nem de médias móveis. Esse resultado é consistente com a
hipótese de eficiência de mercado na forma fraca: os retornos do par BTC/USDT, a despeito
de sua curtose elevada, não apresentam estrutura linear preditível suficiente para justificar
a inclusão de termos AR ou MA. É importante observar que esse resultado não contradiz a
rejeição do Ljung-Box reportada na seção anterior. As autocorrelações detectadas são de
magnitude prática pequena — estatísticas de 50 e 155 para séries com cerca de 200 mil
observações representam efeitos de baixa amplitude — e concentram-se, sobretudo, nos
retornos ao quadrado, não nos níveis; o componente GARCH é o mecanismo mais adequado
para capturá-las.

O resultado mais informativo da seleção refere-se ao componente GARCH. Para as dollar bars,
o modelo GARCH$(1,2)$ --- com um termo de inovação ao quadrado $(\alpha_1)$ e dois termos
de variância defasada $(\beta_1, \beta_2)$ --- foi escolhido em 233 dos 252 folds,
correspondendo a 92,5\% de frequência de seleção. Para as time bars, o modelo GARCH$(1,1)$
canônico foi selecionado na totalidade dos 252 folds (100,0\%). O AIC médio das time bars
ao longo dos folds foi de $-859.083$, contra $-835.475$ para as dollar bars.

A necessidade de um termo $\beta_2$ adicional nas dollar bars tem uma interpretação
econômica plausível. Por construção, as dollar bars possuem espaçamento temporal irregular:
são geradas em maior número durante períodos de alta atividade financeira, quando a
volatilidade do mercado é tipicamente mais elevada. Esse mecanismo faz com que os choques
de volatilidade sejam distribuídos ao longo de mais barras, tornando o decaimento da
variância condicional mais gradual e demandando uma especificação ligeiramente mais rica
para capturá-lo adequadamente. As time bars, ao ignorar esse padrão de irregularidade
temporal, comprimem os choques em barras únicas e produzem uma dinâmica de volatilidade
mais diretamente capturada pelo GARCH$(1,1)$ canônico.

A estabilidade das seleções merece atenção. O fato de o GARCH$(1,1)$ ter sido eleito
em 100\% dos folds para as time bars indica que a dinâmica de volatilidade dessa série
é bem descrita pela especificação padrão da literatura, com pouca ambiguidade quanto à
ordem. Para as dollar bars, a frequência de 92,5\% tampouco é marginal: denota robustez
do GARCH$(1,2)$, não um artefato de um fold específico. Em síntese, o CPCV entregou
seleções estáveis para ambas as séries, com a diferença de especificação entre elas
refletindo uma propriedade genuína das dinâmicas de volatilidade subjacentes.

Cabe destacar que ajuste \textit{in-sample} superior não implica, necessariamente,
melhor desempenho preditivo fora da amostra — distinção fundamental em qualquer
avaliação de modelos de previsão \cite{diebold2015}. Essa questão é endereçada
diretamente na seção seguinte.

A \autoref{tab:modelo_final} apresenta os critérios de informação do
modelo final estimado sobre o conjunto de treino completo.

\begin{table}[H]
\centering
\caption{Critérios de informação do modelo final estimado sobre o
conjunto de treino completo.}
\label{tab:modelo_final}
\begin{tabular}{llSS}
\toprule
{Série} & {Modelo} & {AIC} & {BIC} \\
\midrule
\emph{Dollar Bars} & GARCH(1,2) & -1693794.21 & -1693744.06 \\
\emph{Time Bars}   & GARCH(1,1) & -1742053.06 & -1742012.92 \\
\bottomrule
\end{tabular}
\fonte{Elaboração própria.}
\end{table}

% -----------------------------------------------------------------------------
\section{Desempenho Preditivo Fora da Amostra}
\label{sec:desempenho_previsao}

Esta seção constitui o núcleo analítico do trabalho, respondendo diretamente à pergunta
de pesquisa com base nas métricas calculadas sobre os conjuntos de teste — 171 dias para
as dollar bars (2025-07-14 a 2025-12-31) e 147 dias para as time bars
(2025-08-07 a 2025-12-31) — e comparadas contra o estimador TSRV de volatilidade
realizada como benchmark.

% -----------------------------------------------------------------------------
\subsection{Visualização das Previsões}
\label{subsec:visualizacao_previsao}

Os gráficos a seguir apresentam a série temporal das previsões de variância
condicional agregada ao nível diário em comparação à volatilidade realizada
(TSRV), bem como os erros de previsão ao longo do período de teste.

\begin{figure}[H]
    \centering
    \begin{minipage}[t]{0.48\textwidth}
        \centering
        \includegraphics[width=\textwidth]{analise_previsao/Dollar_Bars_5m/Dollar_Bars_5m_01_serie_temporal_ic.png}
        \subcaption{Série temporal com IC \textit{bootstrap} (95\%).}
        \label{fig:prev_dollar_serie}
    \end{minipage}
    \hfill
    \begin{minipage}[t]{0.48\textwidth}
        \centering
        \includegraphics[width=\textwidth]{analise_previsao/Dollar_Bars_5m/Dollar_Bars_5m_02_erro_previsao.png}
        \subcaption{Erros de previsão.}
        \label{fig:prev_dollar_erro}
    \end{minipage}
    \caption{Previsão da variância condicional versus variância realizada (TSRV) --- \emph{Dollar Bars} GARCH(1,2), período de teste (14/jul/2025--31/dez/2025).}
    \label{fig:prev_dollar_grid}
\end{figure}

\begin{figure}[H]
    \centering
    \begin{minipage}[t]{0.48\textwidth}
        \centering
        \includegraphics[width=\textwidth]{analise_previsao/Time_Bars_5m/Time_Bars_5m_01_serie_temporal_ic.png}
        \subcaption{Série temporal com IC \textit{bootstrap} (95\%).}
        \label{fig:prev_time_serie}
    \end{minipage}
    \hfill
    \begin{minipage}[t]{0.48\textwidth}
        \centering
        \includegraphics[width=\textwidth]{analise_previsao/Time_Bars_5m/Time_Bars_5m_02_erro_previsao.png}
        \subcaption{Erros de previsão.}
        \label{fig:prev_time_erro}
    \end{minipage}
    \caption{Previsão da variância condicional versus variância realizada (TSRV) --- \emph{Time Bars} GARCH(1,1), período de teste (07/ago/2025--31/dez/2025).}
    \label{fig:prev_time_grid}
\end{figure}


% -----------------------------------------------------------------------------
\subsection{Análise de Resíduos}
\label{subsec:analise_residuos}

A análise de resíduos do modelo ajustado constitui um diagnóstico
essencial para verificar se os modelos GARCH estimados capturam
adequadamente a estrutura de dependência das séries de retornos.
Avalia-se, em particular, se os resíduos padronizados apresentam
ausência de autocorrelação serial, distribuição aproximadamente
normal e ausência de efeitos ARCH remanescentes.

% --------------------------------------------------

\begin{figure}[H]
    \centering
    \begin{minipage}[t]{0.48\textwidth}
        \centering
        \includegraphics[width=\textwidth]{analise_residuos/Dollar_Bars_5m/forecast_errors_Dollar_Bars_5m_01_residuos_tempo.png}
        \subcaption{Resíduos padronizados ao longo do tempo.}
        \label{fig:res_dollar_tempo}
    \end{minipage}
    \hfill
    \begin{minipage}[t]{0.48\textwidth}
        \centering
        \includegraphics[width=\textwidth]{analise_residuos/Dollar_Bars_5m/forecast_errors_Dollar_Bars_5m_08_previsao_vs_observado.png}
        \subcaption{Previsão versus variância realizada (TSRV).}
        \label{fig:res_dollar_previsao_obs}
    \end{minipage}
    
    \vspace{0.5cm}
    \begin{minipage}[t]{0.48\textwidth}
        \centering
        \includegraphics[width=\textwidth]{analise_residuos/Dollar_Bars_5m/forecast_errors_Dollar_Bars_5m_04_residuos_vs_ajustados.png}
        \subcaption{Resíduos versus valores ajustados.}
        \label{fig:res_dollar_vs_ajust}
    \end{minipage}
    \hfill
    \begin{minipage}[t]{0.48\textwidth}
        \centering
        \includegraphics[width=\textwidth]{analise_residuos/Dollar_Bars_5m/forecast_errors_Dollar_Bars_5m_02_histograma.png}
        \subcaption{Histograma com curva normal sobreposta.}
        \label{fig:res_dollar_hist}
    \end{minipage}
    \caption{Diagnósticos de ajuste e distribuição dos resíduos --- \emph{Dollar Bars} GARCH(1,2).}
    \label{fig:res_dollar_grid1}
\end{figure}

\begin{figure}[H]
    \centering
    \begin{minipage}[t]{0.48\textwidth}
        \centering
        \includegraphics[width=\textwidth]{analise_residuos/Dollar_Bars_5m/forecast_errors_Dollar_Bars_5m_03_qqplot_normal.png}
        \subcaption{QQ-Plot versus distribuição normal.}
        \label{fig:res_dollar_qq}
    \end{minipage}
    \hfill
    \begin{minipage}[t]{0.48\textwidth}
        \centering
        \includegraphics[width=\textwidth]{analise_residuos/Dollar_Bars_5m/forecast_errors_Dollar_Bars_5m_05_acf_residuos.png}
        \subcaption{Função de Autocorrelação (FAC).}
        \label{fig:res_dollar_acf}
    \end{minipage}
    
    \vspace{0.5cm}
    \begin{minipage}[t]{0.48\textwidth}
        \centering
        \includegraphics[width=\textwidth]{analise_residuos/Dollar_Bars_5m/forecast_errors_Dollar_Bars_5m_06_pacf_residuos.png}
        \subcaption{Função de Autocorrelação Parcial (FACP).}
        \label{fig:res_dollar_pacf}
    \end{minipage}
    \hfill
    \begin{minipage}[t]{0.48\textwidth}
        \centering
        \includegraphics[width=\textwidth]{analise_residuos/Dollar_Bars_5m/forecast_errors_Dollar_Bars_5m_07_acf_residuos_quadrado.png}
        \subcaption{FAC dos resíduos ao quadrado.}
        \label{fig:res_dollar_acf_sq}
    \end{minipage}
    \caption{Diagnósticos de normalidade e dependência --- \emph{Dollar Bars} GARCH(1,2).}
    \label{fig:res_dollar_grid2}
\end{figure}

A análise gráfica dos resíduos padronizados do modelo ajustado sobre as \emph{dollar bars} (Figuras \ref{fig:res_dollar_grid1} e \ref{fig:res_dollar_grid2}) permite avaliar a adequação do GARCH(1,2) sob três aspectos principais: estabilidade temporal, distribuição dos erros e dependência serial.

Com relação à estabilidade temporal e o ajuste geral (\autoref{fig:res_dollar_grid1}), observa-se que os resíduos padronizados (a) distribuem-se de maneira razoavelmente homogênea ao redor de zero, sem apresentar agrupamentos prolongados de alta variância não contabilizada. A dispersão da previsão versus a variância realizada (b) e contra os valores ajustados (c) evidencia a capacidade do modelo em acompanhar a trajetória da volatilidade (TSRV), porém indicando tendência (viés sistemático) em suas previsões. Por fim, o histograma (d) sobreposto à curva normal demonstra como a distribuição empírica dos erros se comporta em relação à distribuição Gaussiana assumida pela Máxima Verossimilhança.

Essa divergência em relação à normalidade é mais claramente constatada nos diagnósticos da \autoref{fig:res_dollar_grid2}. O QQ-Plot (a) revela desvios notáveis nas caudas da distribuição teórica, um indicativo clássico de caudas pesadas (leptocurtose) comum a retornos financeiros, o qual persiste mesmo após a filtragem GARCH. Apesar dos desvios distribucionais nas extremidades, os filtros mostraram-se competentes na remoção da dependência serial: as Funções de Autocorrelação (FAC) e Autocorrelação Parcial (FACP) dos resíduos (b e c) apresentam coeficientes majoritariamente contidos no interior do intervalo de confiança de 95\%, corroborando a adequação da equação da média. Fundamentalmente, a FAC dos resíduos ao quadrado (d) opera como o diagnóstico de efeitos ARCH remanescentes. A ausência de autocorrelações significativas fora das faixas azuis assegura que a estrutura GARCH(1,2) capturou adequadamente a dinâmica de \textit{volatility clustering}, absorvendo a dependência não-linear da série de maneira eficaz.

% --------------------------------------------------

\begin{figure}[H]
    \centering
    \begin{minipage}[t]{0.48\textwidth}
        \centering
        \includegraphics[width=\textwidth]{analise_residuos/Time_Bars_5m/forecast_errors_Time_Bars_5m_01_residuos_tempo.png}
        \subcaption{Resíduos padronizados ao longo do tempo.}
        \label{fig:res_time_tempo}
    \end{minipage}
    \hfill
    \begin{minipage}[t]{0.48\textwidth}
        \centering
        \includegraphics[width=\textwidth]{analise_residuos/Time_Bars_5m/forecast_errors_Time_Bars_5m_08_previsao_vs_observado.png}
        \subcaption{Previsão versus variância realizada (TSRV).}
        \label{fig:res_time_previsao_obs}
    \end{minipage}
    
    \vspace{0.5cm}
    \begin{minipage}[t]{0.48\textwidth}
        \centering
        \includegraphics[width=\textwidth]{analise_residuos/Time_Bars_5m/forecast_errors_Time_Bars_5m_04_residuos_vs_ajustados.png}
        \subcaption{Resíduos versus valores ajustados.}
        \label{fig:res_time_vs_ajust}
    \end{minipage}
    \hfill
    \begin{minipage}[t]{0.48\textwidth}
        \centering
        \includegraphics[width=\textwidth]{analise_residuos/Time_Bars_5m/forecast_errors_Time_Bars_5m_02_histograma.png}
        \subcaption{Histograma com curva normal sobreposta.}
        \label{fig:res_time_hist}
    \end{minipage}
    \caption{Diagnósticos de ajuste e distribuição dos resíduos --- \emph{Time Bars} GARCH(1,1).}
    \label{fig:res_time_grid1}
\end{figure}

\begin{figure}[H]
    \centering
    \begin{minipage}[t]{0.48\textwidth}
        \centering
        \includegraphics[width=\textwidth]{analise_residuos/Time_Bars_5m/forecast_errors_Time_Bars_5m_03_qqplot_normal.png}
        \subcaption{QQ-Plot versus distribuição normal.}
        \label{fig:res_time_qq}
    \end{minipage}
    \hfill
    \begin{minipage}[t]{0.48\textwidth}
        \centering
        \includegraphics[width=\textwidth]{analise_residuos/Time_Bars_5m/forecast_errors_Time_Bars_5m_05_acf_residuos.png}
        \subcaption{Função de Autocorrelação (FAC).}
        \label{fig:res_time_acf}
    \end{minipage}
    
    \vspace{0.5cm}
    \begin{minipage}[t]{0.48\textwidth}
        \centering
        \includegraphics[width=\textwidth]{analise_residuos/Time_Bars_5m/forecast_errors_Time_Bars_5m_06_pacf_residuos.png}
        \subcaption{Função de Autocorrelação Parcial (FACP).}
        \label{fig:res_time_pacf}
    \end{minipage}
    \hfill
    \begin{minipage}[t]{0.48\textwidth}
        \centering
        \includegraphics[width=\textwidth]{analise_residuos/Time_Bars_5m/forecast_errors_Time_Bars_5m_07_acf_residuos_quadrado.png}
        \subcaption{FAC dos resíduos ao quadrado.}
        \label{fig:res_time_acf_sq}
    \end{minipage}
    \caption{Diagnósticos de normalidade e dependência --- \emph{Time Bars} GARCH(1,1).}
    \label{fig:res_time_grid2}
\end{figure}

A avaliação dos resíduos padronizados para o modelo ajustado sobre as \emph{time bars} (Figuras \ref{fig:res_time_grid1} e \ref{fig:res_time_grid2}) segue os mesmos princípios diagnósticos aplicados anteriormente, atestando o desempenho e a adequação da especificação GARCH(1,1) selecionada.

A \autoref{fig:res_time_grid1} ilustra o comportamento geral do ajuste no tempo e suas características distributivas. Os resíduos padronizados (a) exibem homogeneidade ao redor de zero e não demonstram agrupamentos óbvios de alta volatilidade não explicada. A aderência da variância prevista em relação à volatilidade realizada (b) e a distribuição dos resíduos contra os valores ajustados (c) sinalizam a capacidade do modelo em capturar os choques de mercado. O histograma (d), entretanto, mais uma vez evidencia o desafio da premissa Gaussiana: os erros empíricos apresentam um pico excessivo e caudas mais alongadas (leptocurtose) se comparados à curva da distribuição normal esperada, um fenômeno amplamente documentado na modelagem de retornos de alta frequência.

As implicações distributivas são nitidamente corroboradas na \autoref{fig:res_time_grid2}. Os marcantes desvios nas extremidades do QQ-Plot (a) ratificam que a distribuição normal subestima as realizações de cauda (eventos extremos) presentes na série amostrada no tempo. Por outro lado, do ponto de vista de autocorrelação, a filtragem alcança seus objetivos estatísticos: as Funções de Autocorrelação (FAC) e Autocorrelação Parcial (FACP) dos resíduos (b e c) não acusam *lags* significativos persistentes além do intervalo de confiança, sugerindo que qualquer dependência linear foi resolvida. Do mesmo modo, a ausência de correlações relevantes na FAC dos resíduos ao quadrado (d) indica que a estrutura GARCH(1,1) acomodou adequadamente a dependência condicional na série cronológica, eliminando as dinâmicas remanescentes de \textit{volatility clustering}.

% -----------------------------------------------------------------------------
\subsection{Métricas Pontuais}
\label{subsec:metricas_pontuais}

\begin{table}[H]
\centering
\caption{Métricas de avaliação das previsões de variância vs.\ RV
\emph{benchmark}.}
\label{tab:metricas}
\begin{tabular}{lSS}
\toprule
{Métrica} & {\emph{Dollar Bars}} & {\emph{Time Bars}} \\
\midrule
MSE               & 1.05e-7  & 1.65e-7  \\
MAE               & 2.64e-4  & 3.59e-4  \\
QLIKE             & -6.993   & -6.776   \\
$R^2_{\text{MZ}}$ & 0.667    & 0.493    \\
$\hat{\alpha}$    & -9.20e-5 & -2.54e-4 \\
$\hat{\beta}$     & 0.868    & 1.032    \\
\bottomrule
\end{tabular}
\fonte{Elaboração própria.}
\end{table}

As métricas de avaliação pontual, apresentadas na \autoref{tab:metricas},
revelam dominância das dollar bars nas três funções de perda avaliadas. O Erro Quadrático
Médio (MSE) das dollar bars é de $1{,}05 \times 10^{-7}$, contra $1{,}65 \times 10^{-7}$
para as time bars — uma redução de aproximadamente 36\%. O Erro Absoluto Médio (MAE)
exibe comportamento análogo: $2{,}64 \times 10^{-4}$ versus $3{,}59 \times 10^{-4}$,
representando redução de aproximadamente 26\%. A métrica QLIKE, cujos valores menores
(em módulo, mais negativos) indicam melhor desempenho, registra $-6{,}993$ para as dollar
bars e $-6{,}776$ para as time bars, uma diferença de aproximadamente 0,22 em magnitude absoluta.

As três métricas pertencem à classe de funções de perda robustas de \citeonline{patton2011},
cuja propriedade fundamental é preservar o ranking entre previsões de volatilidade mesmo
quando o benchmark utilizado como variável dependente é um proxy ruidoso da volatilidade
latente verdadeira. Essa propriedade é particularmente relevante no presente contexto,
pois o estimador TSRV, ainda que consistente, incorpora ruído de microestrutura e
incerteza de estimação que o afastam da variância quadrática integrada teórica. O fato
de as três métricas robustas apontarem unanimemente na mesma direção reforça a
confiabilidade do resultado.


\subsection{Regressão de Mincer-Zarnowitz}
\label{subsec:mz}

A regressão de Mincer-Zarnowitz \cite{mincer1969}, que regride a volatilidade
realizada diária (TSRV) $\xi_t$ sobre a previsão de variância condicional do modelo
 $\hat{\sigma}^2_t$, fornece um diagnóstico adicional sobre a qualidade e o viés das previsões.
Para as dollar bars, o coeficiente de determinação $R^2$ foi de 0,667 e o intercepto
estimado de $\hat{\alpha} = -9{,}20 \times 10^{-5}$ (p-valor = 0,0168). Para as time bars,
 $R^2 = 0{,}493$ e $\hat{\alpha} = -2{,}54 \times 10^{-4}$ (p-valor = 0,0001).

O modelo baseado em dollar bars explica, portanto, 66,7\% da variação da volatilidade
realizada diária, contra 49,3\% para as time bars. Essa diferença de aproximadamente
17,4 pontos percentuais no $R^2$ é economicamente relevante: para aplicações de gestão
de risco, como o cálculo de VaR (\textit{Value at Risk}) ou a otimização de carteiras
com restrições de variância, modelos que explicam maior parcela da variação da volatilidade
realizada tendem a produzir limites de risco mais calibrados \cite{andersen2003}.

Ambos os interceptos são negativos e estatisticamente significativos (a 5\% para dollar bars e 1\% para time bars), indicando leve
subestimação sistemática da volatilidade realizada por parte de ambos os modelos.
Esse resultado é típico de modelos GARCH paramétricos estimados com inovações gaussianas
quando comparados a estimadores não-paramétricos de alta frequência: o estimador TSRV
captura componentes de variação intrajornada que o GARCH, por ser estimado sobre retornos
diários ou intradiários agregados, não incorpora diretamente. O coeficiente angular $\hat{\beta}$
foi de 0,868 para as dollar bars (p-valor $< 0{,}001$) e de 1,032 para as time bars
(p-valor $< 0{,}001$), sugerindo que o modelo de dollar bars tende a subestimar levemente
a resposta da volatilidade realizada à previsão, enquanto o de time bars apresenta
resposta ligeiramente acima da unidade, indicando leve sobre-reação relativa.

\subsection{Teste de Diebold-Mariano}
\label{subsec:dm}

A comparação formal de capacidade preditiva foi realizada por meio do teste de
Diebold-Mariano \cite{diebold1995}. A estatística DM calculada sobre o subconjunto
de 147 dias em que ambos os modelos produzem previsões sobrepostas foi de $-4{,}806$,
com p-valor bilateral inferior a $0{,}001$.

O sinal negativo da estatística DM indica que o modelo estimado com dollar bars apresenta
perdas (MSE) significativamente menores, em média, do que o modelo ajustado sobre time bars.
O p-valor substancialmente menor que os níveis de significância convencionais (1\%, 5\% e
10\%) permite rejeitar a hipótese nula de igual capacidade preditiva a um nível de
significância de 1\%.

O fato de a rejeição ser estatisticamente forte mesmo em uma amostra de teste
relativamente curta (147 dias de sobreposição) reforça a consistência da superioridade
preditiva das dollar bars evidenciada pelas métricas pontuais. Em termos práticos,
a redução de aproximadamente 36\% no MSE não é um mero artefato amostral, mas uma
diferença de desempenho estrutural e significante que pode se traduzir em melhorias
relevantes para gestão de risco, alocação de carteiras ou precificação de derivativos.

\subsection{Intervalos de Confiança Bootstrap}
\label{subsec:bootstrap}

\begin{table}[H]
\centering
\caption{Intervalos de confiança bootstrap (95\%) para as métricas de
avaliação.}
\label{tab:bootstrap}
\begin{tabular}{llScc}
\toprule
{Série} & {Métrica} & {Pontual} & {IC 95\%} & {SE} \\
\midrule
\multirow{3}{*}{\emph{Dollar Bars}}
    & MSE   & 1.050e-7 & {$[7{,}34\text{e-}8;\; 1{,}49\text{e-}7]$} & 1.93e-8 \\
    & MAE   & 2.637e-4 & {$[2{,}35\text{e-}4;\; 2{,}97\text{e-}4]$} & 1.59e-5 \\
    & QLIKE & -6.993   & {$[-7{,}137;\; -6{,}845]$}                  & 7.48e-2 \\
\addlinespace
\multirow{3}{*}{\emph{Time Bars}}
    & MSE   & 1.645e-7 & {$[1{,}35\text{e-}7;\; 2{,}00\text{e-}7]$} & 1.66e-8 \\
    & MAE   & 3.588e-4 & {$[3{,}29\text{e-}4;\; 3{,}90\text{e-}4]$} & 1.56e-5 \\
    & QLIKE & -6.776   & {$[-6{,}877;\; -6{,}668]$}                  & 5.38e-2 \\
\bottomrule
\end{tabular}
\nota{Valores calculados sobre os 147 dias comuns (07/ago/2025--31/dez/2025).
$B = 2.000$ reamostras; IC pelo método percentil.}
\fonte{Elaboração própria.}
\end{table}

Os intervalos de confiança de 95\% obtidos por bootstrap com $n_{\text{boot}} = 2000$
reamostras, apresentados na \autoref{tab:bootstrap}, complementam a
conclusão do teste DM ao fornecer a margem de incerteza para cada modelo
separadamente.

Para o MSE, os intervalos se sobrepõem ligeiramente: o IC 95\% do MSE das dollar bars é
$[7{,}34 \times 10^{-8},\ 1{,}49 \times 10^{-7}]$, enquanto o das time bars é
$[1{,}35 \times 10^{-7},\ 2{,}00 \times 10^{-7}]$. Essa leve sobreposição não
invalida a conclusão do teste DM, mas reflete a grande variância do MSE, que
penaliza quadraticamente grandes desvios e, portanto, é mais sensível a extremos
em pequenas amostras. Contudo, para o MAE, que é uma métrica absoluta menos
suscetível a \textit{outliers}, os intervalos \textbf{não se sobrepõem}: dollar bars em
$[2{,}34 \times 10^{-4},\ 2{,}97 \times 10^{-4}]$ e time bars em
$[3{,}29 \times 10^{-4},\ 3{,}90 \times 10^{-4}]$. Isso demonstra que a previsão de dollar bars
também possui uma acurácia mediana sistematicamente superior.

O resultado mais robusto do presente trabalho emerge da análise dos intervalos de
confiança para a métrica QLIKE. O IC 95\% das dollar bars é $[-7{,}137,\ -6{,}845]$,
e o das time bars é $[-6{,}877,\ -6{,}668]$. Esses intervalos \textbf{não se sobrepõem},
o que significa que, com 95\% de confiança, a métrica QLIKE das dollar bars é menor do
que a das time bars. A relevância desse achado é amplificada pelo fato de o QLIKE ser,
entre as três funções de perda avaliadas, a mais adequada do ponto de vista teórico para
comparar previsões de volatilidade: conforme demonstrado por \citeonline{patton2011}, é a
única função de perda da família que é robusta ao uso de proxies ruidosos e que ao mesmo
tempo possui propriedades de elicitabilidade para a variância condicional.

Em síntese, o conjunto de evidências — estatística DM negativa e significativa a menos de 1\%,
ICs de MAE e QLIKE não sobrepostos a 95\%, e dominância simultânea em MSE, MAE,
QLIKE e $R^2$ MZ — constitui evidência consistente, formal e coerente em favor das dollar bars
como insumo para modelos GARCH de previsão de volatilidade do BTC/USDT. A superioridade
preditiva documentada é substancial em magnitude e perfeitamente amparada pela estatística.

As \autoref{fig:boot_mse}, \autoref{fig:boot_mae} e \autoref{fig:boot_qlike} apresentam as distribuições
bootstrap das métricas MSE, MAE e QLIKE para cada tipo de barra.

\begin{figure}[H]
    \centering
    \begin{minipage}[t]{0.48\textwidth}
        \centering
        \includegraphics[width=\textwidth]{bootstrap_ic/Dollar_Bars_5m/Dollar_Bars_5m_01_bootstrap_mse.png}
        \subcaption{\emph{Dollar Bars} --- GARCH(1,2).}
    \end{minipage}
    \hfill
    \begin{minipage}[t]{0.48\textwidth}
        \centering
        \includegraphics[width=\textwidth]{bootstrap_ic/Time_Bars_5m/Time_Bars_5m_01_bootstrap_mse.png}
        \subcaption{\emph{Time Bars} --- GARCH(1,1).}
    \end{minipage}
    \caption{Distribuição bootstrap do MSE (2.000 reamostras, IC 95\% pelo método percentil).}
    \label{fig:boot_mse}
\end{figure}

\begin{figure}[H]
    \centering
    \begin{minipage}[t]{0.48\textwidth}
        \centering
        \includegraphics[width=\textwidth]{bootstrap_ic/Dollar_Bars_5m/Dollar_Bars_5m_02_bootstrap_mae.png}
        \subcaption{\emph{Dollar Bars} --- GARCH(1,2).}
    \end{minipage}
    \hfill
    \begin{minipage}[t]{0.48\textwidth}
        \centering
        \includegraphics[width=\textwidth]{bootstrap_ic/Time_Bars_5m/Time_Bars_5m_02_bootstrap_mae.png}
        \subcaption{\emph{Time Bars} --- GARCH(1,1).}
    \end{minipage}
    \caption{Distribuição bootstrap do MAE (2.000 reamostras, IC 95\% pelo método percentil).}
    \label{fig:boot_mae}
\end{figure}

\begin{figure}[H]
    \centering
    \begin{minipage}[t]{0.48\textwidth}
        \centering
        \includegraphics[width=\textwidth]{bootstrap_ic/Dollar_Bars_5m/Dollar_Bars_5m_03_bootstrap_qlike.png}
        \subcaption{\emph{Dollar Bars} --- GARCH(1,2).}
    \end{minipage}
    \hfill
    \begin{minipage}[t]{0.48\textwidth}
        \centering
        \includegraphics[width=\textwidth]{bootstrap_ic/Time_Bars_5m/Time_Bars_5m_03_bootstrap_qlike.png}
        \subcaption{\emph{Time Bars} --- GARCH(1,1).}
    \end{minipage}
    \caption{Distribuição bootstrap do QLIKE (2.000 reamostras, IC 95\% pelo método percentil).}
    \label{fig:boot_qlike}
\end{figure}

% -----------------------------------------------------------------------------
\section{Limitações e Ressalvas}
\label{sec:limitacoes}

A interpretação dos resultados deve ser contextualizada por um conjunto de limitações
metodológicas e de escopo que delimitam o alcance das conclusões.

\subsection{Ativo Único}
\label{subsec:lim_ativo}

Todos os resultados são obtidos para o par BTC/USDT negociado na Binance.
Criptomoedas operam 24 horas por dia, sete dias por semana, e possuem microestrutura
distinta de mercados regulados: ausência de mecanismos de \textit{circuit breaker}, inexistência
de leilão de abertura e fechamento, e liquidez fragmentada entre múltiplas \textit{exchanges}.
Essas características tornam o BTC/USDT um caso de análise com dinâmicas próprias,
e a generalização dos resultados para ações, câmbio ou commodities demanda estudos
adicionais com esses ativos.

\subsection{Período Amostral}
\label{subsec:lim_periodo}

O horizonte de dois anos (2024--2025) é suficiente para estimação e validação dos modelos,
mas pode não capturar toda a variação de regimes de volatilidade relevantes para o Bitcoin.
O período abrange fases de \textit{bull} e \textit{bear market}, mas parâmetros estimados
sobre esse intervalo podem não ser estáveis em regimes futuros com características distintas
das observadas.

\subsection{Calibração do Limiar \texorpdfstring{$\bar{D}$}{D-barra}}
\label{subsec:lim_limiar}

O limiar das dollar bars foi calibrado para produzir o mesmo número médio de barras por dia
que as time bars de 5 minutos, tornando as comparações controladas por volume de observação.
Essa escolha é metodologicamente razoável, mas arbitrária: resultados poderiam diferir com
limiares maiores ou menores, que gerariam séries com granularidade distinta e, possivelmente,
propriedades distribucionais diferentes.

\subsection{Distribuição dos Erros}
\label{subsec:lim_distribuicao}

Ambos os modelos foram estimados com inovações gaussianas. Diante do excesso de curtose
documentado --- especialmente expressivo nas time bars --- distribuições de cauda pesada,
como a t-Student ou a Distribuição de Erro Generalizado (GED), poderiam melhorar o ajuste
e potencialmente alterar a magnitude e o sentido das diferenças entre os modelos.

\subsection{Intervalo de Confiança}
\label{subsec: ic_bootstrap}

Um ponto de atenção sobre os intervalos de confiança, obtidos por meio de bootstrap, é o fato de que a técnica assume que os dados amostrados são independentes e identicamente distribuídos (i.i.d.). Diante da natureza dos dados de série temporal, essa premissa pode não ser totalmente satisfeita, porém observamos na análise nos resíduos de time bars (\autoref{fig:res_time_acf}) e dollar bars (\autoref{fig:res_dollar_acf}) que a autocorrelação é residual, não invalidando completamente o uso do bootstrap.